% Part: set-theory
% Chapter: ordinals
% Section: introduction

\documentclass[../../../include/open-logic-section]{subfiles}

\begin{document}

\olfileid[tr]{sth}{ordinals}{intro}

\olsection{Giriş}

\olref[z][]{chap} içinde, \stagesinf{} ilkesi biçiminde hiyerarşinin
sonsuzuncu bir aşaması bulunduğunu varsaydık (Sonsuzluk Aksiyomumuza da
bakın). Ancak \stagessucc{} verildiğinde sonsuzuncu aşamada duramayız;
ilerlemeyi sürdürmemiz gerekir. Bu nedenle ilk sonsuz aşamadan sonraki
aşamada, ilk sonsuz aşamada kullanılabilir olan kümelerin bütün olası
topluluklarını oluştururuz; sonra bunu tekrarlarız, tekrarlarız,
tekrarlarız, \dots

Burada örtük olarak olan şey, bütün doğal sayılardan \emph{sonra} sayılar
bulunabilmesini sağlayan ``sezgisel'' bir sayı kavramını kullanmaya
başlamamızdır. Söz konusu kavram özellikle \emph{sonlu ötesi sıra sayısı}
kavramıdır. Bu bölümün amacı bu fikri daha kesin hâle getirmektir. Genel sıra
sayısı kavramını inceleyecek, sonra belirli kümeleri açıkça sıra sayılarımız
olarak tanımlayacağız.

\end{document}
